\chapter{Introduction}\label{ch01:chap}

This monograph is a complete, self-contained development of the research
candidate known internally as \emph{Noetic Topos 2.0}: a marked and modal
successor to the conservative Noetic site, together with the finite-context
theory that the successor rests on, the collapse theorem that made a successor
necessary, and the evidence discipline under which every one of these
statements is held. It expands the corresponding v0.2 research paper without
changing any claim, status, or boundary recorded there. The mathematics is
elementary category theory --- categories, presheaves, sieves, Grothendieck
topologies, sheafification --- applied to a bookkeeping problem that arises in
the assessment of AI systems. The bookkeeping problem is real; the mathematics
is standard; the specific construction is a Ring~4 research candidate that has
not been independently reviewed, and this chapter explains, among other
things, exactly what that sentence means and why the monograph repeats it so
often.

The present chapter does six things. It describes the program context in
which the construction arose and the structural (not numerical) difficulty
that motivates it (Sections~\ref{ch01:sec:problem}--\ref{ch01:sec:composition}).
It explains informally why sites and sheaves are the natural candidate
formalism for that difficulty (Section~\ref{ch01:sec:sites}). It recounts the
legacy construction and its intent, and states the collapse problem that the
legacy construction runs into --- informally here, with full proofs deferred
to Chapter~5 (Sections~\ref{ch01:sec:legacy}--\ref{ch01:sec:collapse}). It
presents the three moves that define the successor
(Section~\ref{ch01:sec:threemoves}) and summarizes what the finite theory
actually establishes (Section~\ref{ch01:sec:results}). It gives a detailed
reading guide to the remaining chapters (Section~\ref{ch01:sec:guide}).
Finally, it states the evidence discipline --- rings, ledger proof statuses,
and the precise scope of the executable corroboration --- and fixes the
notation used throughout (Sections~\ref{ch01:sec:evidence}--\ref{ch01:sec:notation}).

\section{The bookkeeping problem}\label{ch01:sec:problem}

AI-assessment work in the Cognitive Architecture System (CAS) program, and in
the Project Coherence effort that grew out of it, increasingly relies on
heterogeneous evidence: observations of a system gathered under different
prompts, different tasks, different roles, and different deployment contexts.
An assessment pipeline may hold, at the same time, transcripts elicited under
one prompting regime, scored behavior under another, structured probe results
from a third, and audit records whose provenance passes through several hands
before reaching the analyst. Each item of evidence is \emph{local}: it was
gathered under specific conditions, and it speaks, in the first instance, only
about the system's behavior under those conditions.

The recurring difficulty is structural rather than numerical. It is not that
the individual observations are noisy --- noise is a familiar problem with
familiar statistical treatments --- but that the observations need not
\emph{compose}. Two local views, each internally coherent, may fail to be
views of any single consistent whole. Three views may be pairwise compatible
and yet jointly incompatible. A conclusion drawn from a family of local views
is only as good as the compatibility of that family, and ordinary aggregation
--- averaging scores, concatenating transcripts, tallying pass rates ---
silently assumes a compatibility that nothing in the pipeline checks.

A concrete schematic (deliberately content-free, in keeping with the claim
discipline of Section~\ref{ch01:sec:evidence}): suppose an assessment records
a system's behavior separately under context $A$ and under context $B$, and
each record is separately well-formed. Whether the two records can be
regarded as restrictions of one description of the system --- whether there
exists a single global section of which both are shadows --- is a further
question, and it is a question with mathematical content. Sometimes the
answer is yes and the gluing is unique; sometimes yes but not uniquely;
sometimes no. A bookkeeping discipline for heterogeneous evidence must be
able to pose this question precisely, answer it where an answer exists, and
refuse to fabricate an answer where none does.

The word ``bookkeeping'' is chosen deliberately. Nothing in this monograph
measures a system, scores a behavior, or validates an assessment method. The
subject matter is the \emph{ledger}: the formal apparatus that records which
local views were gathered, under which admission rule each family of views
was accepted as jointly informative, and what compatibility constraints a
purported global conclusion must satisfy. The apparatus is useful exactly to
the extent that it is honest --- that it refuses to remember less than it
was told and refuses to assert more than it can check --- and the central
technical drama of this monograph, the collapse theorem of Chapter~5, is
precisely a case in which a proposed apparatus turned out, provably, to
remember less than its designers intended.

\section{From aggregation to composition}\label{ch01:sec:composition}

To see why the problem calls for more than statistics, consider the three
questions any discipline for local-to-global reasoning must answer.

\begin{enumerate}
\item \textbf{Which families of local views count?} Not every collection of
observations deserves to be treated as jointly informative about a single
target. A family gathered under systematically related conditions --- for
instance, conditions that jointly exhaust a space of contexts --- has a
claim to joint significance that an arbitrary grab-bag does not. The
discipline must say, for each target of assessment, which families of
``restrictions to sub-contexts'' are admitted as \emph{covers} of that
target.
\item \textbf{When are local data compatible?} Given an admitted family, the
local data attached to its members overlap: two sub-contexts may share a
common further sub-context, and the data restricted to that overlap must
agree if the family is to describe one thing. The discipline must define
compatibility on overlaps and must be able to check it.
\item \textbf{What glues?} Given an admitted family with compatible local
data, the discipline must say whether a global datum exists that restricts
to the given local data, and whether it is unique. Existence without
uniqueness, and uniqueness without existence, are both live possibilities,
and both must be nameable.
\end{enumerate}

Aggregation answers none of these questions; it presupposes affirmative
answers to all three. Composition --- in the mathematician's sense of
building global objects from local ones along specified interfaces --- is
the activity these questions describe, and the mathematics of composition
along covers is sheaf theory. The reader who has met sheaves in geometry
will recognize the three questions as, respectively, the choice of a
Grothendieck topology, the matching-family condition, and the sheaf (gluing)
condition. The reader who has not will find all three developed from
scratch, at tutorial depth, in Chapters~2--4.

\subsection{Desiderata for a bookkeeping discipline}\label{ch01:sec:desiderata}

The three questions can be sharpened into desiderata against which any
candidate formalism --- including the one this monograph develops --- can
be audited. A bookkeeping discipline for heterogeneous assessment evidence
should be:

\begin{enumerate}
\item \emph{Explicit about admission.} The rule by which a family of local
views is admitted as jointly informative must be a registered, inspectable
object, not an implicit habit of the pipeline. (In the formalism: the
topology is declared, not inferred.)
\item \emph{Compositional.} Admitted families must behave coherently under
restriction and iteration: restricting an admitted family to a sub-target
yields an admitted family, and admitting families locally on an admitted
family admits the composite. (The second and third topology axioms.)
\item \emph{Honest about gluing.} The discipline must distinguish, and be
able to certify, the three outcomes: unique gluing, non-unique gluing, and
no gluing. (The sheaf condition as a bijection demand, with its injectivity
and surjectivity halves separately meaningful.)
\item \emph{Provenance-preserving.} The record of \emph{why} each family
was admitted --- under which clause, on which authority, in which review
state --- must survive every derivation the discipline performs. (The
marked-cover layer; by the collapse theorem, this cannot be delegated to
the topology.)
\item \emph{Self-limiting.} The discipline must state the boundaries of
its own guarantees --- which data types it covers, which it provably does
not, and which of its own components are unproven --- with the same
precision as its positive results. (The claim-boundary apparatus of this
monograph.)
\end{enumerate}

Desiderata (1)--(3) are satisfied by the standard site-and-sheaf machinery
and are not project contributions. Desideratum (4) is where the legacy
construction failed and where the successor's principal design work lies.
Desideratum (5) is a discipline on the text itself, enforced throughout by
the evidence rules of Section~\ref{ch01:sec:evidence}.

Two adjacent temptations should be flagged immediately, because the
discipline of this monograph blocks both explicitly (the full treatment is
in the seed paper's gluing-boundary section and recurs in
Chapter~12). First, \emph{probabilistic} local data do not inherit
set-valued gluing by analogy: pairwise-compatible marginal distributions may
admit no global joint distribution, a classical phenomenon studied by
Vorob'ev \cite{vorobev} and, in the sheaf-theoretic language adopted here,
by Abramsky and Brandenburger \cite{abramskybrandenburger}. Second,
\emph{linear} or cellular data require a separate linear representation ---
stalks, inner products, weighted coboundaries --- before spectral tools such
as sheaf Laplacians are even defined \cite{hansenghrist}. Nothing in this
monograph constructs either structure, and no set-valued result proved here
licenses a probabilistic or spectral claim.

\section{Sites and sheaves as the natural candidate}\label{ch01:sec:sites}

Category theory supplies a compact formal answer to question (1) above: a
\emph{site}. A site is a category --- objects are the targets of assessment
and morphisms are the ways one target restricts to, or maps into, another
--- together with a \emph{Grothendieck topology}, an assignment to each
object of a collection of \emph{covering sieves}. A sieve on an object $C$
is a set of morphisms into $C$ closed under precomposition; intuitively, a
downward-closed family of ``ways of looking at $C$ from below.'' The
topology's three axioms (stated and motivated in Chapter~3) guarantee that
the covering notion behaves like a covering notion should: the trivial cover
covers, covers restrict along morphisms, and a family that covers locally on
a cover covers globally.

Questions (2) and (3) are then answered uniformly by the \emph{sheaf
condition}: a presheaf --- an assignment of a set of local data to each
object, contravariant along morphisms --- is a sheaf when, for every
covering sieve, compatible families of local data over the sieve correspond
bijectively to global data over the covered object. The category of sheaves
on a small site is a Grothendieck topos \cite{maclane,johnstone}, and the
inclusion of sheaves into presheaves has a left-exact left adjoint,
\emph{sheafification} \cite{stacks00W1}. Composing sheafification with the
inclusion yields a left-exact idempotent monad on presheaves; this monograph
follows the seed paper in calling such composites \emph{modalities}, and in
calling a modality \emph{trivial} when it is naturally isomorphic to the
identity.

\subsection{A miniature vignette}\label{ch01:sec:vignette}

Because the sheaf condition is the load-bearing idea of the entire
monograph, we first see it informally, on a structure small
enough to hold in the head --- a two-point miniature of the sixteen-context
lattice of Part~II. Let the ``contexts'' be the subsets of a two-element
set $\{1, 2\}$: the empty context $\varnothing$, the two singleton contexts
$\{1\}$ and $\{2\}$, and the full context $\{1,2\}$, ordered by inclusion.
Declare that the family $\{ \{1\} \subseteq \{1,2\},\ \{2\} \subseteq
\{1,2\} \}$ \emph{covers} the full context --- the two singletons jointly
exhaust it --- and that the empty family covers the empty context.

A presheaf of ``local records'' assigns a set of records to each context
and a restriction map to each inclusion. Consider first the presheaf $G$
where a record over a context $V$ is an arbitrary function from $V$ to
$\{0,1\}$ --- a bit written on each point in view --- and restriction
forgets the bits of points no longer in view. Take a compatible pair of
local records over the cover: a bit for point $1$ (a record over $\{1\}$)
and a bit for point $2$ (a record over $\{2\}$); compatibility on the
overlap is automatic, because the overlap is the empty context, over which
$G$ has exactly one record (the empty function). The pair assembles into
exactly one record over $\{1,2\}$, and every record over $\{1,2\}$ arises
this way from exactly one pair. The gluing map is a bijection: $G$ is a
sheaf for this cover. Local bit-records compose perfectly.

Now consider the presheaf $H$ that assigns to \emph{every} nonempty
context the same two-element set $\{\mathsf{p}, \mathsf{q}\}$ --- read: a
global verdict recorded redundantly at each context --- with identity
restrictions, and a single record over the empty context. A compatible
pair over the cover is any pair of verdicts, one over $\{1\}$ and one over
$\{2\}$ (again the empty overlap enforces nothing), so there are four
compatible pairs; but there are only two records over $\{1,2\}$, each of
which restricts to a \emph{matching} pair. The gluing map $2 \to 4$ is
injective but not surjective: the mixed pairs $(\mathsf{p}, \mathsf{q})$
and $(\mathsf{q}, \mathsf{p})$ are locally coherent families that no
global record explains. $H$ is not a sheaf for this cover, and the failure
is informative --- it says that verdict-type data cannot be treated as
freely localizable over these covers, on pain of manufacturing global
verdicts that nothing global supports.

This tiny example already displays every phenomenon the monograph trades
in: covers as admitted families, compatibility on overlaps, gluing as a
bijection demand, sheafhood as a property a presheaf has \emph{relative to
a choice of covers}, and non-sheafhood as a diagnosis rather than a
disaster --- the associated sheaf construction of Chapter~4 repairs $H$
canonically, and what it repairs it to is itself informative. The witness
presheaves that separate the ACBE and RPM modalities in Chapter~9 are
elaborations of exactly the two patterns above, and the reader may find it
useful to keep the pair $(G, H)$ in mind as the germ of that argument.

\subsection{What the formalism promised}\label{ch01:sec:promise}

The promise this formalism held for the assessment program was twofold.
First, the topology is a place to \emph{register admission rules}: to say,
once and in one place, which families of restricted observations count as
covers, and to have the topos machinery enforce the compatibility and gluing
disciplines downstream. Second --- and this was the more ambitious hope ---
different \emph{kinds} of admission rules might generate different
topologies, whose sheafification modalities would then act as distinct
``lenses'' on the same ambient evidence: one modality resolving evidence up
to one kind of cover, another up to another kind, with the interaction of
the lenses itself carrying structure. It is exactly this second promise that
the legacy construction failed to deliver, for reasons that turned out to be
theorems.

\section{The legacy construction and its intent}\label{ch01:sec:legacy}

The conservative Noetic Topos paper --- the ``legacy construction''
throughout this monograph --- introduced a site for the registered category
$\cC$, written $\mathrm{Noetica}$ in the source artifacts: a small concrete
category of finite sets indexed by the forty distinguished CAS indices, with
admissibility-relation generating maps and pullback closure. On $\cC$ it
defined a coverage $J_{\mathrm{legacy}}$ generated by three clauses:

\begin{enumerate}
\item[(O)] \emph{jointly-surjective observational families}: families of
maps into a target whose images jointly exhaust it;
\item[(A)] \emph{designated ACBE descent maps}, including, by reflexivity,
the descent map at equal contexts;
\item[(R)] \emph{the RPM counit} $\varepsilon_X : R(X) \to X$, where the
registered RPM data supply a unit $\eta_X : X \to R(X)$ with
$\varepsilon_X \circ \eta_X = \id_X$ --- a section-retraction pair, and
deliberately \emph{not} an adjunction, in keeping with the program-wide
correction that prerequisite satisfaction is one-directional.
\end{enumerate}

Each clause encodes a defensible bookkeeping intuition, and we record
the intuitions, because the successor preserves them even as it
discards their legacy formalization. Clause (O) says that a family of
observation channels which jointly exhausts a target is entitled to speak
for the target as a whole --- the ordinary local-to-global entitlement.
Clause (A) was meant to encode a descent discipline along designated
context relationships: evidence gathered at related contexts should
descend to, and be answerable at, a common context. Clause (R) was meant
to encode a prerequisite discipline: the counit's passage from a
prerequisite-enriched representation back to the base object was to count
as an evidence-bearing cover, so that satisfaction established at the
enriched level would be creditable at the base. The intuitions survive in
the successor as the two formal coordinates of the context lattice; their
legacy formalizations, as the next section explains, could not.

The legacy paper proved a conditional topos-existence theorem: when $\cC$ is
small, the category of sheaves $\Sh(\cC, J_{\mathrm{legacy}})$ is a
Grothendieck topos. That theorem is valid and is not disputed anywhere in
this monograph. The \emph{intent} attached to the construction, however,
went further: the ACBE clause (A) and the RPM clause (R) were meant to be
remembered by the topology as distinguishable contributions, so that later
work could build one modal operator from the (A)-generated covers and a
different modal operator from the (R)-generated covers --- two separate
lenses on evidence, with the observational clause (O) supplying the common
substrate. The names ACBE and RPM travel through this monograph as formal
labels for these two intended lenses; the reader needs no expansion of the
acronyms, and, as Section~\ref{ch01:sec:evidence} makes precise, the labels
carry no empirical interpretation in any theorem below.

\section{The collapse problem, stated informally}\label{ch01:sec:collapse}

The intent was unrealizable, and the reason is elementary. Two facts about
sieves, each provable in three lines (they are proved, with care, in
Chapter~3), do all the damage:

\begin{enumerate}
\item A sieve containing the identity of its target is the \emph{maximal}
sieve: it contains every morphism into the target, because every morphism
factors through the identity.
\item A sieve containing a \emph{split epimorphism} --- a map with a
one-sided inverse on the right --- contains the identity, by closing under
precomposition with the splitting; hence it is maximal by fact (1).
\end{enumerate}

Now look at the legacy clauses. Clause (A) includes, by reflexivity, the
descent map at equal contexts, which is an identity; every sieve it
generates is maximal by fact (1). Clause (R) includes the counit
$\varepsilon_X$, which is split by the unit $\eta_X$ since
$\varepsilon_X \circ \eta_X = \id_X$; every sieve it generates is maximal by
fact (2). But maximal sieves belong to \emph{every} Grothendieck topology
--- that is the first topology axiom --- so neither clause contributes
anything that the observational clause (O) did not already provide. The
topology generated by all three clauses equals the topology generated by (O)
alone, and any modality built from the (A)- or (R)-generated topology is
trivial. Chapter~5 states this as the Legacy Collapse proposition and proves
it; the registered finite certificate independently re-executes both
mechanisms as the checks
\texttt{legacy\_acbe\_identity\_forces\_maximal\_sieve} and
\texttt{legacy\_rpm\_split\_epi\_forces\_identity}.

Three remarks fix the correct reading of this result.

First, the collapse is \emph{not} a defect of the topos machinery, and it
does not invalidate the legacy topos-existence theorem. What collapses is
the intended interpretive structure. The result is best read as a theorem
about which information a Grothendieck topology is \emph{capable} of
retaining: sieve generation is a closure operator, closure operators forget
how their output was reached, and a topology therefore cannot remember
\emph{why} a family was admitted --- only \emph{that} the sieve it generates
is covering. Provenance is exactly the kind of information a topology
discards by design.

Second, the collapse cannot be escaped by relabeling. Any successor clause
that again contains an identity, or again contains a split epimorphism onto
its target, generates only maximal sieves and collapses the same way. A
successor must \emph{change the generators}, not rename them.

Third, a separate and independent gap accompanies the collapse:
subcanonicity of the legacy site --- the property that representable
presheaves are sheaves, which is what licenses treating a site object as its
own evidence-carrier inside the topos --- is \emph{not derivable} from the
registered axioms. Chapter~6 exhibits a countermodel: a model of the axioms
(finite sets of size at most two, with injections) in which a
jointly-surjective family exists whose gluing map for a representable fails
to be surjective. The consequence is a standing rule, not a repair: until an
exact revised site is proved subcanonical, the distinguished legacy object
is represented in the sheaf topos by its sheafification
$a_{\mathrm{legacy}}(yX)$, never by silently treating $yX$ as a sheaf
\cite{stacks00WO}.

\subsection{Why a topology cannot remember: closure and
provenance}\label{ch01:sec:closure}

We dwell for a moment on the mechanism behind the collapse,
because it explains the shape of the entire successor design and not merely
the failure of one construction. Passing from a generating family of
morphisms to the sieve it generates, and from a collection of generating
sieves to the topology they generate, are both instances of \emph{closure}:
monotone, extensive, idempotent operations that add everything the axioms
force and record nothing about the order or origin of the forcing. Closure
operators are forgetful by design --- that is what makes them canonical.
The closure of a set of premises under logical consequence does not
remember which consequences came from which premises; the subgroup
generated by a set of elements does not remember which generator produced
which product; and the topology generated by a set of covers does not
remember which clause admitted which cover. Asking the generated topology
to distinguish an (A)-cover from an (R)-cover is asking a closure operator
to be injective on histories, which no closure operator is.

Seen this way, the legacy failure decomposes into a general fact and a
specific misfortune. The general fact: even with well-chosen generators,
the topology would remember only the \emph{join} of the admitted covers,
not their labels; any design that needs labels needs a second structure.
The specific misfortune: the legacy generators were not merely unlabeled
but degenerate --- each contained an identity or a split epimorphism, so
each generated only the maximal sieve and contributed nothing at all, not
even unlabeled content. The successor therefore needs two distinct
repairs, and Moves~1 and~2 below are exactly those repairs: new generators
that contribute genuine nonmaximal covers (curing the degeneracy), and an
external provenance record that survives closure (curing the
forgetfulness). Neither repair substitutes for the other. Nonmaximal
covers without external marks would still be anonymous after closure ---
the singleton-context example of Section~\ref{ch01:sec:move2} exhibits
this concretely --- and marks without nonmaximal covers would be labels on
nothing.

\subsection{The two negative results are independent}\label{ch01:sec:independence}

The collapse (Chapter~5) and the subcanonicity gap (Chapter~6) are
logically independent defects, and the monograph keeps them separate
because their repairs are separate. The collapse concerns clauses (A) and
(R) and is repaired by the successor's new generators; it would persist
even if the legacy site were proved subcanonical. The subcanonicity gap
concerns clause (O) --- the jointly-surjective coverage itself --- and is
not repaired by the successor at all: the successor's subcanonicity
theorem (Chapter~9) covers only the new context factor $L$, and the
monograph states explicitly, twice, that it neither repairs nor bears on
the legacy gap. What the successor does about the legacy gap is
disciplinary rather than mathematical: it adopts the representation rule
that bare legacy site objects are sheafified before use whenever
$J_{\mathrm{legacy}}$ is not known subcanonical. A reader auditing the
construction should verify that no argument in Part~II quietly assumes
legacy representables are sheaves; the countermodel of Chapter~6 is the
standing reason such an assumption would be unearned.

\section{The three moves of the successor}\label{ch01:sec:threemoves}

The successor candidate --- Noetic Topos 2.0 --- responds to the collapse
with three moves. Each is developed fully in Part~II of this monograph
(Chapters~7--12); this section states them and their rationale.

\subsection{Move 1: change the generators}\label{ch01:sec:move1}

Since no relabeling of the legacy clauses can produce nonmaximal covers, the
successor introduces a new, separate, finite \emph{context} dimension on
which nontrivial coverage is possible, and leaves the legacy category
untouched. The context factor is the sixteen-element powerset lattice
\[
L = \mathcal{P}(\Omega), \qquad \Omega = \{0,1\} \times \{0,1\},
\]
ordered by inclusion, with $U \to V$ if and only if $U \subseteq V$. The
four atoms of $\Omega$ carry two coordinates; the first coordinate is the
formal ACBE mark and the second the formal RPM mark. Two \emph{transverse
partitions} of $\Omega$,
\[
A_i = \{(i,0),(i,1)\}, \qquad R_j = \{(0,j),(1,j)\}, \qquad i,j \in \{0,1\},
\]
generate two coverages: at a nonempty context $V$, the ACBE generating
family is $\{V \cap A_0 \to V,\ V \cap A_1 \to V\}$ and the RPM generating
family is the analogue with the $R_j$; the empty context is covered by the
null (empty) family. The decisive structural point is local to the global
context: at $\Omega$ the pieces of both generating families are
\emph{proper} subsets ($A_0, A_1 \subsetneq \Omega$ and
$R_0, R_1 \subsetneq \Omega$), so neither global family contains an
identity, and --- a poset having no nontrivial split epimorphisms
(Remark~\ref{ch02:rem:no-split-in-posets}) --- neither contains a split
epimorphism onto its target. The collapse mechanism therefore has nothing
to grip where it mattered. At a context $V$ contained in a single block the
corresponding family degenerates to $\{V \to V,\ \varnothing \to V\}$ and
does generate the maximal sieve; that degeneracy is harmless for the escape
theorem, and it is exactly what makes the external mark record necessary at
singleton contexts (Section~\ref{ch01:sec:move2}). The transversality identity
$A_i \cap R_j = \{(i,j)\}$ --- each ACBE block meets each RPM block in
exactly one atom --- powers essentially every theorem of the finite theory.
Chapter~7 develops the lattice, and Chapter~9 proves that the generated
topologies $J_A$ and $J_R$ are genuine Grothendieck topologies, are
subcanonical, and support distinct, non-identity sheafification modalities.

\subsection{Move 2: externalize provenance}\label{ch01:sec:move2}

The collapse theorem says a topology cannot remember why a cover was
admitted; the successor stops asking it to. Provenance is stored
\emph{outside} the topology, in immutable, hash-chained \emph{marked cover}
records that travel in parallel with the unmarked arrow families the
topology consumes. A marked cover binds, under canonical JSON serialization
and SHA-256 digesting: a derivation kind (generator, pullback, or common
refinement), the target, the exact member arrows, a mark record (cover kind,
source authority reference, assumptions, evidence references, review state,
provenance chain), and the sorted set of parent-cover digests. The integrity
rules are fail-closed: unknown or duplicate keys are rejected, derived
records require a parent resolver that rehashes every parent digest and
recomputes the child, and a common refinement lacking exactly one ACBE and
one RPM generating parent is rejected rather than treated as evidence.

The division of labor is strict and is stated as an invariant of the design:
\emph{the topology consumes the unmarked family; governance and audit code
consume the marked record}. The necessity of the parallel record is not
hypothetical --- it is exhibited inside the finite model itself. At a
singleton context $\{(i,j)\}$, the bare ACBE and RPM families coincide as
arrow lists, so the topology literally cannot distinguish them; the marked
records carry different cover kinds and different digests, and the
certificate confirms both facts
(\texttt{singleton\_unmarked\_families\_coincide},
\texttt{singleton\_marked\_covers\_distinct}). Provenance survives exactly
where the topology forgets. Chapter~8 develops the mechanism, with the
claim boundary --- repeated there and binding here --- that the mark
mechanism is an external serialization and audit invariant, not a proof of
an internal topos object, not a colored-site construction from the
literature, and not yet the subject of a dedicated novelty review.

\subsection{Move 3: name the ambient package honestly}\label{ch01:sec:move3}

The third move is a naming decision with mathematical content. It would be
natural --- and it would be wrong --- to define the successor as the
category of sheaves for the \emph{join} of all the topologies in play, the
``fully resolved'' subtopos. Wrong, because on that subtopos the ACBE and
RPM reflectors provably restrict to the identity (every jointly resolved
object already satisfies both sheaf conditions), which would erase the modal
structure a second time --- an ironic repetition of the collapse by
over-resolution rather than by bad generators. The successor is therefore
defined as the \emph{ambient} package: the presheaf topos
$\mathcal{E}_{\mathrm{amb}} = \PSh(\cC \times L)$, equipped with three
topology refinements (observational in the $\cC$ coordinate, $J_A$ and $J_R$
in the $L$ coordinate), their three sheafification modalities
$M_{\mathrm{obs}}, M_A, M_R$, the resolved subtopos
$\mathcal{E}_{\mathrm{res}}$ for the joined topology, and the external mark
record. The underlying topos candidate is
$\mathrm{NoeTop}_2 = \mathcal{E}_{\mathrm{amb}}$; the full research object
is the structured package
\[
\mathrm{NoeSystem}_2 = \bigl(\mathrm{NoeTop}_2,\ M_{\mathrm{obs}},\ M_A,\
M_R,\ \mathcal{E}_{\mathrm{res}},\ \textsc{external mark record}\bigr),
\]
which is a package, not itself a topos or category, with the provenance
record as external audit data rather than an internal topos object. The
intended discipline: raw and partially resolved evidence lives in the
ambient topos, acted on by the modalities; only jointly resolved objects
enter $\mathcal{E}_{\mathrm{res}}$. Chapter~12 develops the ambient program
and its open obligations, including the explicit boundary that the v0.1
executable schema is deliberately narrower than the ambient construction:
typed observational covers in the $\cC$ coordinate, general arrows of
$\cC \times L$, and full product-site covers remain \textsc{open}.

\section{What the finite theory establishes}\label{ch01:sec:results}

On the context factor $L$ alone --- sixteen objects, small enough to check
exhaustively --- the theory is complete, and Chapters~9--10 prove all of
it. In summary, with ledger identifiers attached as recorded in the
registered theorem ledger:

\begin{enumerate}
\item \emph{Generated topologies} (\textsc{thm-nt2-marked-coverage}): the
ACBE and RPM generating families form pullback-stable coverage bases and
generate Grothendieck topologies $J_A$ and $J_R$ on $L$.
\item \emph{Proper global covers}
(\textsc{thm-nt2-nontrivial-global-covers}): at the global context $\Omega$
both generating families consist of proper pieces, and the generated sieves
are nonmaximal --- the exact point at which the successor escapes the
collapse.
\item \emph{Subcanonicity} (\textsc{thm-nt2-subcanonical-context-factor}):
$J_A$ and $J_R$ are subcanonical on $L$. This concerns $L$ only and neither
repairs nor bears on the legacy subcanonicity gap.
\item \emph{Distinct, non-identity modalities}
(\textsc{thm-nt2-distinct-modalities}): explicit finite witness presheaves
$F_A$ and $F_R$ --- bit assignments on complete blocks --- separate the two
sheaf subcategories: $F_A$ is a $J_A$-sheaf but not a $J_R$-sheaf, and
symmetrically. Hence the reflective subcategories are distinct and neither
sheafification modality is trivial. The scope is $\PSh(L, \Set)$; the
ambient lift proposed in the v0.1 specification contained an error, the
corrected witness remains unverified, and ambient distinctness is
\textsc{open}.
\item \emph{The join} (\textsc{thm-nt2-join-discrete-context}): the common
refinement of the two global covers is the four-singleton cover, and
$J_A \vee J_R$ is the topology generated by the singleton partition ---
equivalently, the ordinary union topology on the finite powerset lattice.
\item \emph{Explicit reflectors} (registered within
\textsc{thm-nt2-modality-composition}; see
Boundary~\ref{ch10:bnd:recap}): a product-over-blocks construction
realizes each associated-sheaf reflector concretely, verified directly
against the universal property rather than through the general plus
construction. This item alone carries no executable corroboration --- the
universal property of the reflector and the sheafhood of its values have no
registered check in either verifier (Chapter~11), and their machine
verification is admission gate~2.
\item \emph{The composition theorem}
(\textsc{thm-nt2-modality-composition}, v0.2): the two modalities satisfy
unit-compatible canonical natural isomorphisms
$M_A M_R \cong M_{\mathrm{atom}} \cong M_R M_A$, and either composite is
naturally isomorphic to sheafification for the joined topology.
\item \emph{Identity on the resolved subcategory}
(\textsc{thm-nt2-resolved-modalities-identity}): every sheaf for the join
is separately a sheaf for each factor topology, so both reflectors restrict
to the identity on the resolved subcategory --- the formal content of
Move~3.
\end{enumerate}

We state one consequence separately, because it corrects a live
narrative risk rather than fulfilling a hope.

\begin{finding}[Distinct but commuting]\label{ch01:finding:commuting}
The finite-context modalities are distinct and non-identity, yet they
commute up to canonical natural isomorphism. Earlier informal program
materials hoped for interacting, order-sensitive modal structure;
distinctness does not imply noncommutation, no current artifact licenses a
noncommutation claim, and any successor construction seeking genuinely
order-sensitive modalities must change the site again rather than
reinterpret this one.
\end{finding}

Beyond the finite theory, the monograph carries exactly one conjecture and a
register of open obligations. The conjecture
(\textsc{thm-nt2-product-site-backward-compatibility}, status
\textsc{proof required}) is the product-site round trip: after atomic
resolution of the context factor, an equivalence
$\mathcal{E}_{\mathrm{res}} \simeq \mathcal{E}_{\mathrm{legacy}}^{4}$ under
which a neutral-channel embedding is fully faithful and evaluation at the
neutral channel is naturally isomorphic to the identity. Two round trips are
required and neither is complete: the categorical comparison, and the strict
serialization requirement that registered identifiers and canonical payload
hashes survive migration bit-for-bit --- a categorical natural isomorphism
is not a serialization receipt. This monograph claims no backward
compatibility. And even a proved round trip would expose an expressivity
boundary, stated in the seed paper and developed in Chapter~12: the resolved
product is four \emph{independent} legacy channels with no remaining
overlap, transition law, or cross-atom coupling, so the sixteen-context
construction is a finite separation and provenance reference model, not a
substantively rich final semantics for contextual cognition. The registered
obligation \textsc{obl-nt2-cross-context-semantics} demands a canonical
decision between naming the ambient modal package as principal semantics and
adding typed coupling with proved coherence; v0.1 licenses neither branch.

\section{Reading guide}\label{ch01:sec:guide}

The monograph is organized in two parts around a hinge. Part~I
(Chapters~2--6) is background and diagnosis: the standard category theory
needed to read everything else, and the two negative results about the
legacy construction. Part~II (Chapters~7--12) is the successor: the finite
context factor, the provenance layer, the finite theorems, the reflector
calculus, the executable certificate, and the ambient program. Four
appendices collect exhaustive finite data. In more detail:

\begin{description}
\item[Chapter 2 (Categories, functors, and presheaves: a primer).] A
tutorial-depth development of categories, functors, natural transformations,
opposite categories, presheaves, representables, the Yoneda lemma (with a
full proof in the poset case, which is the case the finite theory lives in),
and limits with emphasis on pullbacks --- including the fact, used
constantly, that pullbacks in a poset are meets and in a powerset lattice
are intersections. Everything in the chapter is standard material
\cite{awodey,riehl,maclanecwm,maclane}, presented as such, with many worked
examples on two-, three-, and four-element posets chosen to rehearse in
miniature the sixteen-element lattice of Part~II.
\item[Chapter 3 (Sieves, Grothendieck topologies, and sites).] Sieves and
their identification with subfunctors of representables; pullback sieves;
the three topology axioms with intuition; coverage bases and the generated
topology; the two elementary maximality facts that drive the collapse; a
complete enumeration of all Grothendieck topologies on the two-element
chain; and the union coverage on a powerset lattice. Again standard
material \cite{maclane,stacks00ZC,johnstone}, at tutorial depth.
\item[Chapter 4 (Sheaves, sheafification, and modalities).] Matching
families, the sheaf condition, sheafification and its exactness
\cite{stacks00W1}, subcanonicity and representable sheaves
\cite{stacks00WO}, and the definition of modality (left-exact idempotent
monad from sheafification-then-inclusion) with the triviality criterion.
\item[Chapter 5 (The legacy construction and the collapse).] The registered
legacy site in full; the Legacy Collapse proposition with proof; what the
collapse does and does not invalidate; the corrigendum reading.
\item[Chapter 6 (The countermodel).] Non-derivability of legacy
subcanonicity: the injections-on-small-sets countermodel, worked in
complete detail, and the resulting representation rule for legacy objects.
\item[Chapter 7 (The context lattice).] The lattice $L$, its
combinatorics (sixteen contexts, eighty-one inclusion arrows), the
transverse partitions, the generating families, and the transversality
identity.
\item[Chapter 8 (Marked covers).] The provenance layer: record schemas,
canonical serialization, hash chaining, derivation kinds, the fail-closed
integrity rules, the singleton-context necessity example, and the claim
boundary on the mechanism's mathematical status.
\item[Chapter 9 (The finite theorems).] Full proofs of the coverage-basis
lemma, the generated-topology theorem, properness of global covers,
subcanonicity, the witness-presheaf distinctness theorem, and the join
identification.
\item[Chapter 10 (Reflectors and composition).] The explicit
product-over-blocks reflector; its sheafhood; the universal-property
verification; the composition theorem with its unit-compatibility and
naturality checks; identity on the resolved subcategory; the
distinct-but-commuting finding.
\item[Chapter 11 (The executable certificate).] What the finite-site
checker and the modality-composition verifier actually execute, item by
item; the formal (Lean) layer and its limits; and the certificate's own
claim boundary, quoted and adopted.
\item[Chapter 12 (The ambient program).] The ambient package
$\mathrm{NoeSystem}_2$; the gluing boundary for probabilistic and linear
data; the backward-compatibility conjecture and the expressivity boundary;
related work and novelty status; the full register of open obligations and
admission gates.
\item[Appendix A (Context tables).] The sixteen contexts of $L$ with their
partition pieces and inclusion data.
\item[Appendix B (Sieve census).] The enumeration underlying the
certificate's sieve counts.
\item[Appendix C (Witness tables).] The witness presheaves $F_A$ and $F_R$
tabulated on all sixteen contexts, with the gluing computations displayed.
\item[Appendix D (Glossary).] Project terms, ledger identifiers, and status
vocabulary.
\end{description}

The logical dependencies among the chapters are summarized below; a
chapter presupposes the chapters listed against it and nothing else beyond
Chapter~1.

\begin{center}
\begin{tabular}{lll}
\toprule
Chapter & Depends on & Character \\
\midrule
2 & --- & standard background \\
3 & 2 & standard background \\
4 & 2, 3 & standard background \\
5 & 3 & project diagnosis (negative) \\
6 & 3, 4 & project diagnosis (negative) \\
7 & 2, 3 & project construction \\
8 & 7 & project construction (external layer) \\
9 & 4, 7 & project theorems (finite, proved) \\
10 & 9 & project theorems (finite, proved) \\
11 & 9, 10 & corroboration record \\
12 & 5--10 & program, conjecture, obligations \\
A--D & as cited & exhaustive finite data \\
\bottomrule
\end{tabular}
\end{center}

\emph{Prerequisites and pedagogy.} The monograph presupposes comfort with
naive set theory --- sets, functions, images, unions, intersections,
partitions --- and with reading and writing elementary proofs; a reader
wanting a refresher on proof technique may consult a standard
introduction such as Hammack's \emph{Book of Proof} \cite{bookofproof}. No
prior category theory is assumed: Chapters~2--4 build everything used
later, and they build it on posets wherever posets suffice, because the
finite theory of Part~II lives entirely on a poset and the poset case
strips away every distraction. The tutorial chapters follow three habits
throughout. First, every abstract definition is followed immediately by
its instantiation on a two-, three-, or four-element poset, and the same
small posets recur so that the reader accumulates a portable laboratory.
Second, enumerations are actually carried out --- all sieves on a small
object, all topologies on the two-element chain --- rather than gestured
at, because the exhaustive habit is precisely the habit the certificate of
Chapter~11 mechanizes at scale. Third, exercises are interleaved where
working a computation by hand is more instructive than reading it;
solutions are not provided, but every exercise is within reach of the
worked examples that precede it, and none is needed for the logical
development.

\emph{Suggested paths.} A reader new to category theory should read
Chapters~2--4 in order, then 5--6, then Part~II linearly; the small-poset
examples of Chapters~2--3 are designed so that every later construction has
already been met in miniature. A reader fluent in topos theory can skim
Chapters~2--4 for notation, read Chapter~5 and Chapter~6 for the negative
results (both are short), and then work through Chapters~7--12, consulting
the appendices as the exhaustive data are cited. A reader whose interest is
the evidence discipline itself --- what is proved, what is corroborated,
what is open --- should read Section~\ref{ch01:sec:evidence} below,
Chapter~11, and Chapter~12, in that order, before anything else.

\section{Evidence discipline}\label{ch01:sec:evidence}

Nothing in this monograph is peer-reviewed, and the monograph is written so
that this fact is impossible to lose. This section states the discipline
once, in full; individual chapters repeat the locally relevant boundaries in
\emph{Claim Boundary} environments at the point of use.

\subsection{Artifact status and rings}\label{ch01:sec:rings}

The governing artifact status, inherited verbatim from the v0.2 seed paper
and its source basis, is: \textsc{draft / gate 1 research candidate /
ring 4 / not licensed / held}. Ring~4 is the program's outermost evidence
ring: material that may be developed, executed, and internally corroborated,
but that licenses no external claim --- no empirical claim, no causal claim,
no safety claim, no product claim --- and that may leave the ring only
through the explicit admission gates listed in Chapter~12. Those gates
include, among others: independent category-theory review of every
proposition and theorem; proof-assistant formalization of the finite site,
topologies, witnesses, reflectors, join, and natural isomorphisms; the
typed observational-cover construction; the product-site lift; the
backward-compatibility conjecture with both round trips; the
cross-context-semantics decision; separate probability-gluing theorems; a
dedicated novelty search for the marked-cover mechanism; registry,
authorship, and IP reconciliation; and external referee review with
explicit confirmation that no empirical or product wording was licensed.
Drafting this monograph promotes no claim and closes no gate.

\subsection{Source basis}\label{ch01:sec:sources}

Every project fact, status, hash, and count in this monograph comes from the
registered source artifacts and from nowhere else:
\begin{itemize}
\item the formal specification \texttt{NOETIC\_TOPOS\_2\_SPEC\_v0.1.md}, with
SHA-256 digest
\begin{center}
\texttt{156A5D8C5320F36E81E3B46473BA8A038B35672DCD28839A404D96A3DB1A763A};
\end{center}
\item the theorem ledger v0.1 (digest abbreviated \texttt{52E00065}\ldots\ in
the seed paper, full value in the registered ledger);
\item the proof-obligation ledger v0.1 (abbreviated \texttt{FD567E38}\ldots);
\item the finite model certificate
\texttt{FINITE\_MODEL\_CERTIFICATE\_v0.1.json}, model identifier
\textsc{nt2-context-2x2-transverse-partitions};
\item the composition proof
\texttt{MODALITY\_COMPOSITION\_PROOF\_v0.1.md} and its certificate;
\item the publication delta at base commit\\
\texttt{a1cd990109e7561dfd03949bf0e1653461c7ae44};
\item the reference implementation under \texttt{src/noetic\_topos2/}.
\end{itemize}
Where this monograph
states a number --- sixteen contexts, eighty-one arrows, and so on --- the
number is either proved in the text from the definitions or attributed to
the certificate; where a detail is unspecified in the sources, the monograph
says so rather than supplying one.

\subsection{Ledger proof statuses}\label{ch01:sec:statuses}

Every theorem-shaped statement in this monograph carries the proof status
recorded in the registered theorem ledger, and the status vocabulary is
fixed as follows. \textsc{informal proof present} means a written
mathematical proof exists in the source artifacts and is reproduced or
expanded here; all finite-context theorems of Chapters~9--10 carry this
status, together with exhaustive executable corroboration.
\textsc{conjecture} / \textsc{proof required} means precisely what it says;
the backward-compatibility statement is the only such item.
\textsc{open} marks constructions specified but not built or not proved
(the typed observational covers, the ambient lift of distinctness, the
product-site lift of the composition theorem, the cross-context semantics
decision). \textsc{novelty unassessed} is carried by all eight
theorem-ledger records: no dedicated literature search has been performed,
and this monograph accordingly never calls any project construction a new
theorem in the scholarly sense. Standard results are cited to the standard
literature \cite{maclane,johnstone,stacks00ZC,stacks00W1,stacks00WO} and
presented as standard; the boundary between textbook material and project
candidate is marked in every chapter.

\subsection{What exhaustive finite corroboration does and does not
mean}\label{ch01:sec:corroboration}

The finite theory of Part~II is small enough to check by exhaustive
computation, and the project does so with two independent executable
layers, described fully in Chapter~11. The finite-site checker verifies,
over all four atoms and sixteen contexts: properness and joint
exhaustiveness of the global covers, pullback stability of every generating
cover, sheafhood of all sixteen representables, functoriality and the four
sheaf/non-sheaf verdicts for the witness presheaves, the transverse
refinement identity, the singleton coincidence-and-distinctness pair for
marked covers, marked-cover round-trip and parent-provenance binding, and
both legacy collapse mechanisms --- twenty-three registered checks, all
passing. The modality-composition verifier computes exact finite
Grothendieck topologies by enumerating all 298 sieves over the sixteen
contexts, saturating maximality, pullback stability, and local character to
a fixed point; it reports 168 covering sieves for both the joined and
singleton presentations, verifies the Grothendieck axioms independently on
both computed topologies, and checks the composite decompositions, both
canonical natural isomorphisms, inverse laws, restriction naturality,
commutation naturality, and unit compatibility across 18 executable
fixtures and 3{,}012 restriction squares. At the bound commit, the
finite-site suite (8 tests), the modality-composition suite (6 tests), and
the pinned Lean Gate~1 verifier pass; the Lean scope covers four-atom
exhaustiveness, proper global pieces, transverse refinement, and
witness-map lemmas, and does \emph{not} formalize the sixteen presheaves,
sheaf conditions, topology generation, reflector universal properties, the
composition theorem, the product-site comparison, or the legacy round trip.

The two executable layers differ in method: the finite-site checker
(described in the sources as dependency-free) works from the coverage bases
and checks basis-level conditions directly, while the composition verifier
constructs the full topologies by saturating all sieves under maximality,
pullback stability, and local character to a fixed point and then
re-verifies the Grothendieck axioms independently on the computed result.
The methodological difference means the two layers do not simply re-run one
encoding of ``topology'' twice; the deterministic run at the bound commit
passes both.

What this corroboration \emph{means}: for a statement quantified over a
finite structure, an exhaustive check of every instance is a genuine and
strong form of evidence --- there is no sampling, no approximation, and no
untested case within the model checked. It caught, during development,
exactly the kinds of slips (a nonfunctorial witness candidate, an unstated
empty-context convention) that informal proofs are prone to.

What it does \emph{not} mean, in the certificate's own words, adopted here
as governing: exhaustive execution over one sixteen-object site
``corroborates the finite proofs; it is not proof-assistant verification,
external peer review, empirical validation, or a product-claim license.''
Corroboration is relative to the implementation's encoding of the
definitions; a shared misreading of a definition by prover and programmer
would survive it. That is why proof-assistant formalization and independent
review remain listed as open admission gates rather than as formalities.

\subsection{Formal marks carry no empirical
interpretation}\label{ch01:sec:marks-boundary}

\begin{boundary}\label{ch01:bdy:marks}
The coordinate names ACBE and RPM, wherever they appear in Part~II, are
formal labels preserved for controlled migration vocabulary. They carry no
empirical interpretation until a separate operationalization and validation
record exists. No theorem in this monograph attaches meaning to them beyond
their combinatorics, and no theorem here establishes that either coordinate
detects a cognitive, causal, safety, or alignment property.
\end{boundary}

\subsection{How to read a claim in this monograph}\label{ch01:sec:reading-claims}

Because the status vocabulary carries real weight, two worked parses may
help calibrate the reader.

\emph{Parse 1: a finite theorem.} The distinctness theorem of Chapter~9
(ledger \textsc{thm-nt2-distinct-modalities}) asserts that two explicitly
constructed presheaves separate the two sheaf subcategories of
$\PSh(L, \Set)$. Its full epistemic position is: a written informal proof
exists (reproduced and expanded in Chapter~9); every step quantifies over
a sixteen-object structure and has been exhaustively executed by the
certificate (the checks \texttt{acbe\_witness\_functorial},
\texttt{acbe\_witness\_is\_acbe\_sheaf} true,
\texttt{acbe\_witness\_is\_rpm\_sheaf} false, and their RPM-side
counterparts); the statement's scope is $\PSh(L, \Set)$ and \emph{not} the
ambient category, because the v0.1 ambient-lift witness was found to be
erroneous and the corrected ambient statement is \textsc{open}; the ledger
record includes the ambient-lift sentence, and the publication layer ---
followed by this monograph --- corrects it; the result is
\textsc{novelty unassessed} and has no independent review. Every clause of
that position is part of what the theorem \emph{is}, here.

\emph{Parse 2: the conjecture.} The backward-compatibility statement
(ledger \textsc{thm-nt2-product-site-backward-compatibility}) is written
in theorem-like language --- an equivalence, a fully faithful embedding, a
natural isomorphism --- but carries status \textsc{proof required}. In
this monograph such a statement licenses exactly nothing: no chapter may
use it as a premise, no summary may report it as a property of the
construction, and the correct citation form is ``conjectured, proof
required.'' The same reading applies to every \textsc{open} item: the
ambient lift, the typed observational covers, the product-site lift of the
composition theorem, and the cross-context-semantics decision are
specifications of future work, not results.

The general rule the parses instantiate: in this monograph, a claim is the
conjunction of its mathematical statement and its recorded status, and
quoting the former without the latter is a misquotation.

\subsection{Relation to the seed artifacts}\label{ch01:sec:seed-relation}

The monograph's textual relationship to its sources is expansion, not
revision. The v0.2 research paper (which itself superseded a v0.1 white
paper) fixed the claim set: the collapse diagnosis, the countermodel, the
finite-context theorems and composition law, the marked-cover mechanism,
the ambient package, the conjecture, and the open-obligation register.
This monograph adds tutorial background (Chapters~2--4), fully expanded
proofs where the paper compressed them, worked small-poset laboratories,
exhaustive appendix tables, and extended discussion of the evidence
discipline --- and it adds no new mathematical claim, promotes no status,
and closes no obligation. Where the monograph's expanded prose and a
source artifact could be read to differ, the source artifact governs; the
statuses, hashes, and counts of Section~\ref{ch01:sec:sources} are the
fixed points of that ordering. Readers who need the compressed,
single-artifact form of the material should read the v0.2 paper; readers
who need the executable ground truth should read the certificate and the
implementation under \texttt{src/noetic\_topos2/}; this monograph sits
between them as the didactic layer.

\subsection{Writing rules}\label{ch01:sec:rules}

\begin{ruleenv}[Claim discipline for this monograph]\label{ch01:rule:claims}
(i) Standard mathematics is presented at tutorial depth, cited to the
standard literature, and never presented as project novelty. (ii) Project
constructions are presented with their ledger statuses attached and are
never called novel; all carry \textsc{novelty unassessed}. (iii) No
sentence asserts an empirical, causal, safety, or product property of any
AI system. (iv) Hashes, counts, and identifiers appear only as recorded in
the source artifacts; unspecified details are declared unspecified. (v)
Claim boundaries travel with the material they bound: each chapter restates
the boundaries relevant to its content.
\end{ruleenv}

\section{Notation and conventions}\label{ch01:sec:notation}

\begin{notation}\label{ch01:not:general}
Categories are written in calligraphic or roman letters; $\cC$ is reserved
for the registered legacy category ($\mathrm{Noetica}$ in the source
artifacts). $\Set$ is the category of sets and functions in a fixed
Grothendieck universe; smallness assumptions are stated where used.
$\Ob(\cC)$ is the class of objects; $\Hom_{\cC}(A,B)$ the set of morphisms
from $A$ to $B$; $\id_A$ the identity. Composition is written
$g \circ f$ for $f : A \to B$ followed by $g : B \to C$, applied right to
left. $\PSh(\cC)$ denotes the category of presheaves
$\cC^{\mathrm{op}} \to \Set$, and $\Sh(\cC, J)$ the full subcategory of
sheaves for a topology $J$. The representable presheaf of an object $C$ is
written $yC = \Hom_{\cC}(-, C)$.
\end{notation}

\begin{notation}\label{ch01:not:posets}
A poset $(P, \le)$ is regarded as a category with a unique morphism
$U \to V$ exactly when $U \le V$; Chapter~2 verifies that this is a
category and that all diagrams in it commute. In the powerset lattice
$\mathcal{P}(X)$ the order is inclusion, meets are intersections, and joins
are unions. The successor's context lattice is
$L = \mathcal{P}(\Omega)$ with $\Omega = \{0,1\} \times \{0,1\}$; its
atoms are written $(i,j)$, its ACBE blocks $A_i$, its RPM blocks $R_j$,
and its generated topologies $J_A$, $J_R$, with join
$J_A \vee J_R$. The seed artifacts' internal label $J_{\mathrm{atom}}$ for
the join is retained where the ledger is quoted, with the standing caution
that it does not assert the standard category-theoretic atomic topology;
$J_{\mathrm{singleton}}$ is the preferred publication-layer name.
\end{notation}

\begin{notation}\label{ch01:not:project}
Modalities are composites $M = i a$ of sheafification $a$ followed by
inclusion $i$; $M_{\mathrm{obs}}$, $M_A$, $M_R$ are the observational,
ACBE, and RPM modalities of the ambient construction. The ambient and
resolved categories are $\mathcal{E}_{\mathrm{amb}} = \PSh(\cC \times L)$
and $\mathcal{E}_{\mathrm{res}} = \Sh(\cC \times L, J_{\mathrm{res}})$
with $J_{\mathrm{res}} = J_{\mathrm{obs}} \vee J_A \vee J_R$;
$\mathcal{E}_{\mathrm{legacy}} = \Sh(\cC, J_{\mathrm{legacy}})$. Ledger
identifiers are set in small capitals (\textsc{thm-nt2-marked-coverage});
executable check names and file names are set in typewriter face with
escaped underscores (\texttt{finite\_site.py}). Statuses are set in small
capitals: \textsc{open}, \textsc{proof required}, \textsc{informal proof
present}, \textsc{novelty unassessed}, \textsc{not licensed}.
\end{notation}

\begin{notation}[Environment conventions]\label{ch01:not:environments}
The monograph uses a fixed repertoire of numbered environments, and their
names are part of the claim discipline. \emph{Theorem}, \emph{Proposition},
\emph{Lemma}, and \emph{Corollary} are reserved for statements with written
proofs in this text; when such a statement corresponds to a ledger record,
the ledger identifier is displayed with the statement. \emph{Conjecture}
marks the backward-compatibility statement and nothing else.
\emph{Definition}, \emph{Construction}, \emph{Example}, and
\emph{Exercise} have their usual meanings; \emph{Notation} fixes symbols.
\emph{Remark} is expository and carries no claim weight. \emph{Claim
Boundary} (the \texttt{boundary} environment) is normative: it restricts
what the surrounding material may be taken to assert, and it binds every
reading of that material. \emph{Finding} records a corrective result ---
something the mathematics settled against a prior informal expectation ---
and \emph{Rule} records a standing discipline adopted by the text.
\emph{Design argument} (the \texttt{designargument} environment) is an
unnumbered block that draws a design or policy consequence from an
already-proved result; it is deliberately not a proof, asserts no new
mathematics, and carries no tombstone. \emph{Axiom}, \emph{Algorithm},
\emph{Protocol}, \emph{Invariant}, \emph{Check Class}, and \emph{Threat} are
numbered alongside the theorem-class environments and name what they say.
Proofs end with the usual tombstone; proof sketches are labeled as
sketches and are used only for standard results whose full proofs are
cited.
\end{notation}

\begin{notation}\label{ch01:not:typography}
The monograph source is ASCII; all mathematical symbols are produced by
LaTeX commands. Numbered environments (theorem, proposition, lemma,
definition, example, remark, boundary, and the rest) are shared across the
monograph with per-chapter numbering; labels are prefixed by chapter slug.
Tables use booktabs rules. Displayed equations are numbered only when
referenced.
\end{notation}

The following summary table collects the recurring symbols for reference.

\begin{center}
\begin{tabular}{ll}
\toprule
Symbol & Meaning \\
\midrule
$\cC$ & registered legacy category ($\mathrm{Noetica}$) \\
$J_{\mathrm{legacy}}$ & legacy coverage, clauses (O), (A), (R) \\
$\Omega$ & the four-atom set $\{0,1\} \times \{0,1\}$ \\
$L$ & context lattice $\mathcal{P}(\Omega)$, sixteen objects \\
$A_i$, $R_j$ & ACBE and RPM partition blocks \\
$J_A$, $J_R$ & generated context topologies \\
$J_A \vee J_R$ & joined (singleton) topology \\
$F_A$, $F_R$ & witness presheaves for distinctness \\
$M_A$, $M_R$, $M_{\mathrm{obs}}$ & sheafification modalities \\
$\mathcal{E}_{\mathrm{amb}}$, $\mathcal{E}_{\mathrm{res}}$ & ambient and resolved categories \\
$\mathrm{NoeTop}_2$, $\mathrm{NoeSystem}_2$ & topos candidate; full research package \\
$yC$ & representable presheaf of $C$ \\
$m_C$ & maximal sieve on $C$ \\
\bottomrule
\end{tabular}
\end{center}

\begin{remark}\label{ch01:rem:tone}
A last word on tone. The monograph's subject is a repair, and the repair is
presented together with the failure it repairs, the certificate that
corroborates it, and the list of everything it does not yet establish. This
is deliberate. In the program's experience, the value of a bookkeeping
formalism is inseparable from the honesty of its own bookkeeping: a
construction that overstated its status would fail the very discipline it
proposes. Preserving the boundaries with the same precision as the theorems
is part of the result.
\end{remark}
